% ============================================================
%  Глава 8 — MPC: горизонт, лифтинг, QP
% ============================================================
\chapter{MPC: Model Predictive Control}
\label{ch:mpc}

\section{Вопрос}

ЛКР минимизирует $J$ на \emph{бесконечном} горизонте, выдаёт постоянное
усиление $\mat{K}$ и не знает о ограничениях $\vect{u}_\text{min} \le
\vect{u} \le \vect{u}_\text{max}$. Эти два недостатка устраняет MPC.

\begin{ideabox}
\textbf{Идея MPC:} вместо одного решения на $\infty$ --- решать
\emph{маленькую} оптимизационную задачу \emph{на каждом тике}, на
конечный горизонт $N$ вперёд, с явными ограничениями.

\textbf{Receding horizon (скользящий горизонт):}
\begin{enumerate}
  \item Сейчас, в момент $k$, посмотри на $N$ шагов вперёд.
  \item Сосчитай \emph{весь будущий план} $\vect{u}_0, \vect{u}_1, \ldots,
        \vect{u}_{N-1}$, минимизируя стоимость.
  \item Примени \emph{только первый шаг} $\vect{u}_0$ к роботу.
  \item На следующем тике начни заново с $\vect{x}_{k+1}$.
\end{enumerate}
Это похоже на шахматы: считаем дерево на N ходов, делаем один, считаем
заново.
\end{ideabox}

\section{Постановка}

\begin{formulabox}[title={MPC-задача (с ограничениями)}]
\begin{equation}
  \begin{aligned}
  \min_{\,\vect{u}_0, \ldots, \vect{u}_{N-1}}\;
  & J = \sum_{i=0}^{N-1}\Big(\vect{x}_i\T \mat{Q}\vect{x}_i +
                              \vect{u}_i\T \mat{R}\vect{u}_i\Big)
        + \vect{x}_N\T \mat{P}_f\,\vect{x}_N\\
  \text{при условиях } &
        \vect{x}_{i+1} = \mat{A}_d\vect{x}_i + \mat{B}_d\vect{u}_i,
        \quad i = 0, \ldots, N-1,\\
  & \vect{u}_\text{min} \le \vect{u}_i \le \vect{u}_\text{max},
        \quad i = 0, \ldots, N-1,\\
  & \vect{x}_0 \text{ задано (измерено сейчас)}.
  \end{aligned}
  \label{eq:mpc_problem}
\end{equation}
\end{formulabox}

Новое:
\begin{itemize}
  \item Сумма по \emph{конечному} $N$ (не до бесконечности).
  \item \textbf{Терминальный штраф} $\mat{P}_f$ --- штраф за последнее
        состояние $\vect{x}_N$. Если выбрать $\mat{P}_f$ как решение DARE
        для $(\mat{A}_d, \mat{B}_d, \mat{Q}, \mat{R})$, то конечный горизонт
        ведёт себя «как бесконечный» --- сумма от $N$ до $\infty$
        заменяется аккуратно квадратичным выражением.
  \item Явные box-ограничения на $\vect{u}_i$.
\end{itemize}

\section{Лифтированная динамика}

Переменными оптимизации в \eqref{eq:mpc_problem} являются $\vect{u}_0, \ldots,
\vect{u}_{N-1}$ --- всего $rN$ скаляров. Состояние $\vect{x}_i$ не
независимая переменная: оно \emph{зависит} от $\vect{u}_i$ через
$\vect{x}_{i+1} = \mat{A}_d\vect{x}_i + \mat{B}_d\vect{u}_i$. Удобно это
зависимое представление сделать явным.

\begin{ideabox}
\textbf{Лифтинг (lifting).} Разворачиваем рекурсию по динамике.
\begin{align*}
  \vect{x}_1 &= \mat{A}_d\vect{x}_0 + \mat{B}_d\vect{u}_0,\\
  \vect{x}_2 &= \mat{A}_d\vect{x}_1 + \mat{B}_d\vect{u}_1
              = \mat{A}_d^2\vect{x}_0 + \mat{A}_d\mat{B}_d\vect{u}_0 + \mat{B}_d\vect{u}_1,\\
  \vect{x}_3 &= \mat{A}_d^3\vect{x}_0 + \mat{A}_d^2\mat{B}_d\vect{u}_0
              + \mat{A}_d\mat{B}_d\vect{u}_1 + \mat{B}_d\vect{u}_2,\\
             & \vdots
\end{align*}
Складываем все состояния в большой столбец $\vect{X}$ и все управления
--- в $\vect{U}$:
\end{ideabox}

Введём блочные векторы:
\[
  \vect{X} = \begin{pmatrix} \vect{x}_1\\ \vect{x}_2\\ \vdots\\ \vect{x}_N \end{pmatrix} \in \Real^{nN},
  \qquad
  \vect{U} = \begin{pmatrix} \vect{u}_0\\ \vect{u}_1\\ \vdots\\ \vect{u}_{N-1} \end{pmatrix} \in \Real^{rN}.
\]

Тогда вся динамика «помещается» в одно матричное равенство:

\begin{formulabox}[title={Лифтированная динамика}]
\begin{equation}
  \vect{X} \;=\; \mat{\Phi}\,\vect{x}_0 \;+\; \mat{\Gamma}\,\vect{U},
  \label{eq:lifted}
\end{equation}
где
\[
  \mat{\Phi} = \begin{pmatrix}
    \mat{A}_d\\ \mat{A}_d^2\\ \vdots\\ \mat{A}_d^N
  \end{pmatrix} \in \Real^{nN \times n},
  \qquad
  \mat{\Gamma} = \begin{pmatrix}
    \mat{B}_d & 0 & \cdots & 0\\
    \mat{A}_d\mat{B}_d & \mat{B}_d & \cdots & 0\\
    \vdots & & \ddots & \vdots\\
    \mat{A}_d^{N-1}\mat{B}_d & \mat{A}_d^{N-2}\mat{B}_d & \cdots & \mat{B}_d
  \end{pmatrix} \in \Real^{nN \times rN}.
\]
\end{formulabox}

\begin{codebox}[title={\filepath{pi\_nodes/control/mpc\_controller.py:134}}]
\begin{lstlisting}[language=Python]
def _build_qp_matrices(self):
    n, r, N = self.n, self.r, self.N
    Phi = np.zeros((n * N, n))
    Gamma = np.zeros((n * N, r * N))
    Apow = np.eye(n)
    for i in range(N):
        Apow = self.Ad @ Apow                    # A^{i+1}
        Phi[i*n:(i+1)*n, :] = Apow
        for j in range(i + 1):
            block = np.linalg.matrix_power(self.Ad, i - j) @ self.Bd
            Gamma[i*n:(i+1)*n, j*r:(j+1)*r] = block
    # ... дальше Q_bar, R_bar, H, f_mat
\end{lstlisting}
\end{codebox}

\subsection{Стоимость в лифтированной форме}

Соберём блочно-диагональные веса (последний блок --- терминальный):

\[
  \bar{\mat{Q}} = \begin{pmatrix}
    \mat{Q} &       &        & \\
            & \mat{Q} &      & \\
            &       & \ddots & \\
            &       &        & \mat{P}_f
  \end{pmatrix} \in \Real^{nN\times nN},
  \qquad
  \bar{\mat{R}} = \mat{I}_N \otimes \mat{R} \in \Real^{rN\times rN}.
\]

Подставим \eqref{eq:lifted} в стоимость $J$:

\[
  J(\vect{U}) = (\mat{\Phi}\vect{x}_0 + \mat{\Gamma}\vect{U})\T \bar{\mat{Q}}
                 (\mat{\Phi}\vect{x}_0 + \mat{\Gamma}\vect{U})
               + \vect{U}\T \bar{\mat{R}}\,\vect{U}.
\]

Раскрываем и группируем по степеням $\vect{U}$:

\begin{formulabox}[title={Квадратичная форма стоимости}]
\begin{equation}
  J(\vect{U}) = \half\,\vect{U}\T \mat{H}\,\vect{U} + \vect{f}(\vect{x}_0)\T \vect{U} + \text{const},
  \label{eq:quad_cost}
\end{equation}
\[
  \mat{H} = 2\,(\mat{\Gamma}\T \bar{\mat{Q}}\mat{\Gamma} + \bar{\mat{R}}),
  \qquad
  \vect{f}(\vect{x}_0) = 2\,\mat{\Gamma}\T \bar{\mat{Q}}\mat{\Phi}\,\vect{x}_0.
\]
\end{formulabox}

«const» --- слагаемое, зависящее только от $\vect{x}_0$. Для оптимизации
по $\vect{U}$ оно неважно.

\section{Квадратичное программирование (QP)}

С учётом ограничений MPC превращается в каноническую \textbf{QP-задачу}:

\begin{formulabox}[title={QP-постановка MPC}]
\begin{equation}
  \min_{\vect{U}}\; \half \vect{U}\T \mat{H}\vect{U} + \vect{f}\T \vect{U},
  \qquad
  \vect{U}_\text{min} \le \vect{U} \le \vect{U}_\text{max}.
  \label{eq:qp}
\end{equation}
\end{formulabox}

Здесь $\vect{U}_\text{min}, \vect{U}_\text{max}$ --- повторение
$\vect{u}_\text{min}, \vect{u}_\text{max}$ по $N$ блокам.

\textbf{Свойства задачи:}
\begin{itemize}
  \item $\mat{H} = \mat{H}\T$, положительно определённая
        $\Rightarrow$ задача \textbf{выпуклая}, имеет единственный
        минимум.
  \item Ограничения --- линейные неравенства, выпуклые.
  \item Существуют десятки QP-солверов (OSQP, qpOASES, scipy.optimize и др.).
\end{itemize}

\section{Два режима в проекте: \texttt{clip} и \texttt{qp}}

В \filepath{pi\_nodes/control/mpc\_controller.py} реализованы \textbf{две}
стратегии решения:

\subsection{\texttt{clip} (по умолчанию)}

\textbf{Без} ограничений (\emph{в задаче, не в железе}) решение QP
\eqref{eq:qp} аналитическое --- производная по $\vect{U}$, равная нулю:

\[
  \mat{H}\,\vect{U}^* + \vect{f} = 0
  \;\;\Rightarrow\;\;
  \vect{U}^* = -\mat{H}^{-1}\vect{f}(\vect{x}_0).
\]

Берём первые $r$ компонент --- это $\vect{u}_0$. После лишних выкладок
получается, что

\begin{equation}
  \vect{u}_0 = -\mat{K}_\text{first}\,\vect{x}_0,
  \qquad
  \mat{K}_\text{first} = (\mat{H}^{-1}\,2\mat{\Gamma}\T \bar{\mat{Q}}\mat{\Phi})_{[1\ldots r,\,:]}.
  \label{eq:K_first}
\end{equation}

То есть аналитический MPC --- это \textbf{та же линейная обратная связь},
что ЛКР, только усиление $\mat{K}_\text{first}$ другое (зависит от $N$).
После умножения мы \emph{клиппим} управление к ограничениям:

\[
  \vect{u}_0 \leftarrow \operatorname{clip}(\vect{u}_0, \vect{u}_\text{min}, \vect{u}_\text{max}).
\]

\begin{notebox}
\textbf{Когда это оптимально:} ровно когда невозмущённое решение
$-\mat{K}_\text{first}\vect{x}_0$ \emph{и так} попадает в box. Тогда
клип ничего не меняет. Если выходит за границы --- клиппинг даёт
\emph{допустимое}, но \emph{не оптимальное} управление. В большинстве
рабочих режимов робота этого хватает.
\end{notebox}

\begin{codebox}[title={\filepath{mpc\_controller.py:288}}]
\begin{lstlisting}[language=Python]
def step(self, x, x_ref=None):
    if x_ref is None:
        dx = x
    else:
        dx = x - x_ref
    if self._solver == "qp":
        return self._solve_qp(dx)
    # Explicit clip
    u = -self.K_first @ dx
    return np.clip(u, self.u_min, self.u_max)
\end{lstlisting}
\end{codebox}

\subsection{\texttt{qp} (полный)}

При \texttt{control.mpc.solver: qp} в конфиге проект решает QP \emph{честно}:

\begin{codebox}[title={\filepath{mpc\_controller.py:303}}]
\begin{lstlisting}[language=Python]
def _solve_qp(self, dx):
    f = self.f_mat @ dx
    H = self.H
    u_lo = np.tile(self.u_min, self.N)
    u_hi = np.tile(self.u_max, self.N)
    bounds = list(zip(u_lo, u_hi))
    U0 = -np.linalg.solve(H, f)              # warm start
    U0 = np.clip(U0, u_lo, u_hi)
    res = minimize(lambda U: 0.5*U@H@U + f@U,
                   U0, jac=lambda U: H@U + f,
                   method="L-BFGS-B", bounds=bounds)
    return (res.x if res.success else U0)[:self.r]
\end{lstlisting}
\end{codebox}

Используется L-BFGS-B (scipy) с warm start от аналитического решения.
Делает то же, что OSQP, только медленнее. \emph{Включают только если
clip регулярно упирается в ограничения}.

\section{Выбор горизонта $N$}

\begin{ideabox}
\textbf{Большой $N$} (например, 50):
\begin{itemize}
  \item Регулятор «видит дальше» --- может заранее тормозить, чтобы
        не перелететь цель.
  \item QP становится больше: размер $\vect{U}$ растёт как $rN$,
        размер $\mat{H}$ --- как $(rN)^2$. Время на тик растёт нелинейно.
\end{itemize}

\textbf{Маленький $N$} (например, 1):
\begin{itemize}
  \item Решается быстро.
  \item Регулятор «близорукий», поведение ближе к ЛКР с малым весом
        конечного штрафа.
\end{itemize}

\textbf{Рабочий компромисс:} $N \cdot T_s \gtrsim 2 \tau_\text{slow}$, где
$\tau_\text{slow}$ --- самая медленная постоянная времени объекта. Для
нас $\tau_\text{slow} = \tau_v = 0.15$~с, $T_s = 0.02$~с $\Rightarrow$
$N \ge 15$. В дефолтах --- $N = 10$ или $20$ (см. \texttt{mps.horizon\_N}).
\end{ideabox}

\section{Терминальный штраф $\mat{P}_f$}

\begin{ideabox}
\textbf{Проблема конечного горизонта:} что мешает MPC сделать «хитрое»
движение, которое в окне $[0, N]$ выглядит выгодно, а после $N$
оставляет робот «в плохой позе» (далеко от цели)? Ничего, если стоимость
после $N$ мы игнорируем.

\textbf{Решение:} добавить терминальный штраф $\vect{x}_N\T \mat{P}_f\,
\vect{x}_N$, который аппроксимирует «стоимость до бесконечности» из
$\vect{x}_N$.

Лучшая такая аппроксимация --- решение DARE \eqref{eq:dare} для
$(\mat{A}_d, \mat{B}_d, \mat{Q}, \mat{R})$, то есть $\mat{P}_f = \mat{P}_\text{LQR}$.
Тогда MPC \emph{в пределе} $N \to \infty$ сходится к ЛКР, а на конечном
$N$ ведёт себя «как ЛКР, но с уважением к ограничениям».
\end{ideabox}

\begin{codebox}[title={\filepath{mpc\_controller.py:126}}]
\begin{lstlisting}[language=Python]
def _compute_terminal_penalty(self):
    """If Pf not given, solve DARE for a guaranteed-stable terminal cost."""
    return solve_discrete_are(self.Ad, self.Bd, self.Q, self.R)
\end{lstlisting}
\end{codebox}

\section{Слежение за reference}

До сих пор минимизировали $\vect{x}\T \mat{Q}\vect{x}$ --- штраф за
отклонение от \emph{нуля}. Но цель сценария --- $\vect{x}^\text{ref} \ne 0$
(например, $s_\text{ref}$ растёт линейно к $D$). Стандартный приём ---
работать в \emph{ошибке слежения} $\Delta\vect{x} = \vect{x} - \vect{x}^\text{ref}$:

\begin{equation}
  J = \sum_{i=0}^{N-1}\Big((\vect{x}_i - \vect{x}^\text{ref})\T \mat{Q}\,(\vect{x}_i - \vect{x}^\text{ref})
        + \vect{u}_i\T \mat{R}\vect{u}_i\Big)
        + (\vect{x}_N - \vect{x}^\text{ref})\T \mat{P}_f\,(\vect{x}_N - \vect{x}^\text{ref}).
  \label{eq:mpc_tracking}
\end{equation}

В коде --- именно так:

\begin{codebox}[title={\filepath{mpc\_controller.py:288}}]
\begin{lstlisting}[language=Python]
def step(self, x, x_ref=None):
    dx = x - x_ref  if x_ref is not None  else x
    ...
\end{lstlisting}
\end{codebox}

\section{Sanity-check замкнутой системы}

После \texttt{rebuild()} (Apply из UI) можно убедиться, что новый
регулятор устойчив:

\begin{codebox}[title={\filepath{mpc\_controller.py:332}}]
\begin{lstlisting}[language=Python]
def closed_loop_eigenvalues(self):
    """|λ| < 1 ⇔ closed-loop x_{k+1} = (Ad - Bd K_first) x is stable."""
    return np.linalg.eigvals(self.Ad - self.Bd @ self.K_first)

def is_stable(self):
    return bool(np.all(np.abs(self.closed_loop_eigenvalues()) < 1.0 - 1e-9))
\end{lstlisting}
\end{codebox}

Этот тест в UI делает «Validate matrices» --- кнопка перед Apply.

\section{Hot-reload между прогонами}

В МПС-учебном модуле пользователь \emph{редактирует} матрицы из UI и
жмёт Apply. На Pi прилетает \texttt{mps/matrices/set}, и нода вызывает
\texttt{mpc.rebuild()}: атомарно пересчитываются $\mat{\Phi}, \mat{\Gamma},
\mat{H}, \mat{K}_\text{first}$, при ошибке --- rollback к предыдущему
снимку. Это и есть «учебный пульт».

\begin{codebox}[title={\filepath{mpc\_controller.py:172}, фрагмент rebuild}]
\begin{lstlisting}[language=Python]
def rebuild(self, *, Ad=None, Bd=None, Q_diag=None, R_diag=None,
            N=None, u_min=None, u_max=None, Pf=None):
    snapshot = (self.Ad.copy(), ..., self.K_first.copy())
    try:
        # обновить поля, провалидировать формы, пересчитать matrices
        self._build_qp_matrices()
    except Exception:
        # Rollback on ANY failure
        (self.Ad, ..., self.K_first) = snapshot
        raise
\end{lstlisting}
\end{codebox}

\section{Что мы получили}

\begin{itemize}
  \item \textbf{MPC} --- решение QP-задачи на каждом тике с явными
        ограничениями.
  \item Через \emph{лифтинг} динамики QP сводится к канонической форме
        $\half\vect{U}\T \mat{H}\vect{U} + \vect{f}\T \vect{U}$.
  \item Без ограничений решение аналитическое, $\vect{u}_0 =
        -\mat{K}_\text{first}\vect{x}_0$. С ограничениями --- настоящая
        QP-задача, в проекте --- L-BFGS-B с warm start.
  \item Дефолтный режим --- \texttt{clip}: считаем $\mat{K}_\text{first}$,
        умножаем, клипуем. Быстро и в большинстве случаев оптимально.
  \item Терминальный штраф $\mat{P}_f$ --- решение DARE; делает короткий
        горизонт «бесконечно-устойчивым».
  \item Горизонт $N$ --- компромисс: $N T_s \gtrsim 2 \tau_\text{slow}$.
  \item \texttt{rebuild()} --- атомарное пересоздание контроллера с
        rollback на ошибку; используется при Apply из UI.
\end{itemize}

Теперь, наконец, всё готово, чтобы запустить настоящий сценарий.
